\documentclass[11pt]{article}
\usepackage[margin=1in,top=0.9in,bottom=0.9in]{geometry}
\usepackage{amsmath,amssymb,amsthm,hyperref}
\hypersetup{colorlinks=true,linkcolor=blue,urlcolor=blue,citecolor=blue}
\setlength{\emergencystretch}{2em}
\newtheorem{theorem}{Theorem}
\title{An exact counterexample to trigonometric positivity\\for square roots of geometric sums}
\author{Research proof and formalization packet}
\date{October 26, 2026}
\begin{document}
\maketitle
\begin{abstract}
We disprove Kovalev's Conjecture~1.3 on positive trigonometric partial
sums of the Taylor coefficients of $\sqrt{1+z+\cdots+z^d}$.
An exact counterexample has $d=49$, partial-sum index $949$, and
$t=\arccos(523/527)$. A finite integer certificate proves the sum is
strictly negative. The normalized square-root definition, coefficient
recurrence, integer certificate, cosine identity and final source
counterexample are all checked by Lean. The separate zero-free
Taylor-polynomial conjecture is not refuted by this example.
\end{abstract}

\section{The source conjecture and counterexample}
For $d\geq2$, define real coefficients $b_k^{(d)}$ by the branch at the
origin with value one:
\[
 \sqrt{1+z+\cdots+z^d}=\sum_{k\geq0}b_k^{(d)}z^k.
\]
Kovalev's Conjecture~1.3 \cite{K}, indexed as MathDB~369172, states that
\[
 \mathcal T_n^{(d)}(t)=\sum_{k=0}^n b_k^{(d)}\cos(kt)>0
 \qquad(d\geq2,\ n\geq0,\ t\in\mathbb R).
\]
The statement is present both in arXiv:2507.04544v1 and in the published
2026 paper, page~274.

\begin{theorem}\label{thm:counterexample}
With $t_0=\arccos(523/527)$,
\[
 \mathcal T_{949}^{(49)}(t_0)<0.
\]
\end{theorem}

Its approximate value is $-0.000475321859370024$; the proof below uses
exact integers, not rounding. This is one counterexample outcome. The
index, rational angle, auxiliary recurrences and certificate receive no
separate discovery credits. Conjecture~1.2 in the same source concerns
zero-free Taylor polynomials. Negating a sufficient positivity condition
does not negate that different conjecture, and no such claim is made.

\section{Identifying the source coefficients exactly}
Write $D=d+1$, $G_D(z)=1+z+\cdots+z^{D-1}$, and let $F_D$ be its
normalized formal square root. It exists, for example, by substituting
$G_D-1$ in the formal binomial series with exponent $1/2$.
The identities $F_D(0)=1$ and $F_D^2=G_D$ also select it uniquely:
two such roots have $(F_1-F_2)(F_1+F_2)=0$ in the integral domain of
formal series; the second factor has constant term two and cannot vanish.
Thus these are exactly the Taylor coefficients specified in the source.

Since $F_D^2(1-z)=1-z^D$, formal differentiation and multiplication give
\begin{equation}\label{eq:ode}
 2(1-z)(1-z^D)F_D'
   =\bigl(1+(D-1)z^D-Dz^{D-1}\bigr)F_D.
\end{equation}
Writing $F_D=\sum b_kz^k$ and comparing the coefficient of $z^{k-1}$
yields, for $k\geq1$,
\begin{align}\label{eq:rec}
 2k b_k={}&(2k-1)b_{k-1}
  +\mathbf1_{k\geq D}(2k-3D)b_{k-D}\notag\\
 &-\mathbf1_{k\geq D+1}(2k-3D-1)b_{k-D-1}.
\end{align}
Terms with negative indices are omitted. Since $2k\ne0$, the initial
value $b_0=1$ and this recurrence determine every needed coefficient.

For $D=50$, the packet supplies integer data $B_0,\ldots,B_{949}$ and
checks $B_0=1$ and, for $1\leq k\leq949$,
\begin{align}\label{eq:intrec}
 2k B_k={}&4(2k-1)B_{k-1}
   +\mathbf1_{k\geq50}4^{50}(2k-150)B_{k-50}\notag\\
 &-\mathbf1_{k\geq51}4^{51}(2k-151)B_{k-51}.
\end{align}
By induction using \eqref{eq:rec}, these finite integer checks prove
$B_k=4^k b_k^{(49)}$ throughout the required range. No guessed
coefficient values or unproved recurrence assumptions enter the final
source theorem.

\section{A finite integer sign certificate}
Let $p=523$, $q=527$, and define integers
\[
 C_0=1,\quad C_1=523,\quad
 C_k=1046C_{k-1}-277729C_{k-2}\quad(k\geq2).
\]
The cosine addition identity and $\cos t_0=p/q$ give, by induction,
\[
 C_k=q^k\cos(kt_0).
\]
Our finite $C$ table is checked against these exact initial values and
recurrences through index~949.

Finally put
\[
 H_0=0,\qquad H_{k+1}=2108H_k+B_kC_k\quad(0\leq k\leq949).
\]
Induction in Horner's recurrence gives
\begin{equation}\label{eq:horner}
 H_{950}=2108^{949}\sum_{k=0}^{949}
       \frac{B_k}{4^k}\frac{C_k}{527^k}
       =2108^{949}\mathcal T_{949}^{(49)}(t_0).
\end{equation}
The certificate computes $H_{950}<0$ in exact integer arithmetic. Its
absolute value has 3152 decimal digits; the complete integer and
denominator $2108^{949}$ are included in the packet rather than printed
as several pages of digits. Since the multiplier in \eqref{eq:horner}
is positive, this proves Theorem~\ref{thm:counterexample}.

For reproducibility, the finite data can be generated from
\[
 (1-z^{50})^{1/2}(1-z)^{-1/2}
  =\left(\sum_{j\geq0}a_jz^{50j}\right)
   \left(\sum_{r\geq0}\frac{\binom{2r}{r}}{4^r}z^r\right),
\]
where $a_0=1$ and $a_j=a_{j-1}(2j-3)/(2j)$.
The generator uses only rational arithmetic. Its correctness is not a
trusted premise: Lean independently checks \eqref{eq:intrec}, the $C$
recurrences and the negative Horner result for the literal integer data.

\section{Formal verification and prior-art review}
The seven-module Lean certificate has 55 inventoried declarations.
The source root is constructed unconditionally from Mathlib's formal
binomial series. The source-to-recurrence bridge and finite recurrence
uniqueness identify its coefficients with the certified data. The final
endpoint \texttt{KovalevCounterexample.explicit\_source\_counterexample}
has no assumptions. The companion endpoint formally negates the universal
positivity conjecture. No analytic remainder bound, floating-point sign
decision or missing source hypothesis is assumed.

The integer checks use ordinary kernel-checked \texttt{decide}; no
\texttt{native\_decide}, added axiom or proof hole is used. The complete
certificate replays in pinned Lean~4.34.1. The only reported axioms are
\texttt{propext}, \texttt{Classical.choice}, and \texttt{Quot.sound}.
The tables are encoded by balanced integer lookups to keep replay
practical; this affects representation, not the statement or trust model.

Independent exact computations by direct square-root convolution and
scaled integer or quadratic-ring cosine arithmetic agree with the
negative value. A separate reviewer checked the mathematical and source
bridges, including $d=49$ versus $D=50$ and the $H_{950}$ indexing.
The actual public publisher PDF was inspected and retains Conjecture~1.3.
Documented literature searches located no prior counterexample to this
exact statement. This is a finite novelty investigation, not an absolute
certification of worldwide priority. No minimality in $d$ or $n$ is claimed.

\begin{thebibliography}{9}
\bibitem{K} L. V. Kovalev,
\emph{Sharp bounds for some segments of bounded power series},
Arkiv f\"or Matematik 64(2) (2026), 273--282, Conjecture~1.3.
\href{https://doi.org/10.4310/ARKIV.2026.v64.n2.a4}{DOI: 10.4310/ARKIV.2026.v64.n2.a4}.
Preprint: \href{https://arxiv.org/abs/2507.04544v1}{arXiv:2507.04544v1} (2025).
\end{thebibliography}
\end{document}
